% TOP_PROBLEM: {"id": 30001953, "title": "Hopf Conjecture", "category_id": 6, "category": {"id": 6, "name": "geometry", "display_name": "Geometry", "description": "Euclidean and non-Euclidean geometry, geometric structures.", "slug": "geometry", "order_index": 6, "created_at": "2026-07-31T15:26:25.671Z"}, "status": "open"}
% FIRST_POSTED: 2026-09-25
% ENTRY_KIND: statement_only
% !TeX program = lualatex
% Definition and sources only; this file is not an LLM solution attempt.
\documentclass[11pt]{article}
\usepackage[margin=25mm]{geometry}
\usepackage{fontspec}
\setmainfont{Latin Modern Roman}
\usepackage{amsmath,amssymb,mathtools}
\usepackage{unicode-math}
\setmathfont{Latin Modern Math}
\usepackage{xurl,hyperref}
\hypersetup{colorlinks=true,urlcolor=blue}
\setcounter{secnumdepth}{0}
\setlength{\parindent}{0pt}
\setlength{\parskip}{5pt}
\setlength{\emergencystretch}{3em}
\begin{document}
\section*{\texorpdfstring{Hopf Conjecture}{Hopf Conjecture}}\label{30001953}
\subsection{Definitions and mathematical statement}
Let $(M^{2m},g)$ be a closed Riemannian manifold. For a two-plane $\sigma=\operatorname{span}(u,v)\subset T_xM$, sectional curvature is
$K_g(\sigma)=\langle R(u,v)v,u\rangle/(|u|^2|v|^2-\langle u,v\rangle^2)$, with the convention making the round sphere positive. The Euler characteristic is
$\chi(M)=\sum_j(-1)^j\dim H_j(M,{\mathbb{Q}})$. The conjecture states
\[
 K_g(\sigma)>0\text{ for every }x,\sigma
 \quad\Longrightarrow\quad\chi(M)>0.
\]
No symmetry or simple-connectedness hypothesis is imposed. Positivity of Ricci curvature or scalar curvature alone is a different and weaker curvature condition.
\par\smallskip\textit{Formulation and concept references:} \href{https://arxiv.org/abs/2507.16936}{[S1]}.

\subsection{Short English statement}
Must every closed even-dimensional manifold with everywhere positive sectional curvature have positive Euler characteristic?
\subsection{Sources}
\begin{itemize}
\item[S1]Lee Kennard, Lawrence Mouille, Jan Nienhaus. On Hopf's conjecture and positive second intermediate Ricci curvature. 2025.
\url{https://arxiv.org/abs/2507.16936}.
\textit{Formulation source.}
\end{itemize}
\end{document}
